\renewcommand*{\authorOfNextLines}{HTL Weiz}
\label{par:KurzfassungDiplomarbeitErlaeuterung}
Die Kurzfassung soll eine interessierte Öffentlichkeit in die Lage versetzen, die durch die Diplomarbeit erzielte Leistung, 
 insbesondere die ingenieurmäßige Eigenleistung der Verfasserinnen \bzw der Verfasser und die damit nachgewiesenen Kompetenzen zu erkennen und einzuschätzen. 
Um diesen Zweck zu erreichen, ist eine Gleiderung der Kurzfassung in die Abschnitte
\begin{itemize}
	\item Aufgabenstellung
	\item Realisierung
	\item Ergebnisse
	\item Allgemeine Angaben
\end{itemize} 
vorgesehen.\\
Im Folgenden werden für jeden Abschnitt typische Fragen formuliert, auf die die Kurzfassung eine Antwort geben soll. 
Die Struktur mit den typischen Fragen ist als Leitlinie zu verstehen, die auf den Großteil der Diplomarbeiten angewendet werden kann.\\
Wie die Diplomarbeit ist auch der Text der Kurzfassung in wissenschaftlicher Form darzustellen (keine \gqq{ICH/WIR-Sätze}).

